\section{Probing the Saturated Concept: A Valence Direction in LLMs}
\label{sec:vaxis-llms}

To test that the saturation regularity operates on a real semantic
axis rather than a noise dimension, we extract one explicitly. The
construction uses nine sentences. Given a language model, we author
one short emotion-evocative story per FACED class and generate $50$
paraphrases per story by independent LLM calls. For each class we
average the language model's last-token activation (at the
penultimate layer) over all $50$ paraphrases; the resulting vector is
the class \emph{centroid}---a single point in feature space that
summarises the model's representation of that emotion. Principal
component analysis (PCA) on the nine centroids gives the V-axis as
the first principal component, oriented so the Joy centroid projects
positive. We use Qwen3-4B as the lead model and verify the recipe is
robust to paraphrase budget, prompt rephrasing, and model choice
(Appendix~\ref{app:vaxis-protocol}).

\paragraph{The same direction across modalities.}
The same nine-story V-axis recipe, applied in each modality's native
encoder, recovers human affect across three settings that share no
training data: language sentiment, brain responses, and visual scenes
(Table~\ref{tab:cross-modal}). The text result is a true zero-shot
projection; the EEG result uses the same recipe re-extracted from
CLIP-text plus a ridge linear probe (LOOCV) on $160$ EEG features;
the vision result uses CLIP-image of nine matched images.

\begin{table}[h]
\centering
\small
\setlength{\tabcolsep}{6pt}
\begin{tabular}{@{}lllrl@{}}
\toprule
Modality & Axis source & Probe & Score & Reference \\
\midrule
Text   & Qwen3-4B (9 stories) & none (proj.) & $\mathbf{0.832}$ AUC & matches sup.\ LR on $5$k examples \\
Brain  & CLIP-text (9 stories) & ridge LOOCV & $\mathbf{r{=}0.87}$ & $p<10^{-9}$, see \S\ref{sec:vaxis-brain} \\
Vision & CLIP-image (9 images) & none (proj.) & $\mathbf{r{=}0.869}$ & beats sup.\ ridge \\
\bottomrule
\end{tabular}
\caption{Cross-modal external validity of the V-axis \emph{recipe}
(PCA on nine emotion-class centroids). The axis source differs per
row (Qwen3-4B for SST-2, CLIP-text for cohort EEG, CLIP-image for
OASIS); the Qwen3-4B V-axis projected onto cohort EEG via the same
ridge gives $r{=}0.80$ (\S\ref{sec:vaxis-brain}). Full sentiment
(IMDB, TweetEval, Rotten Tomatoes), lexicon (Hu \& Liu, Warriner),
and multilingual results in Appendix~\ref{app:per-llm-deep-dive}.}
\label{tab:cross-modal}
\end{table}

The LLM-side direction is a strong sentiment classifier: $0.832$
AUC on SST-2 matches a supervised logistic regression on $5$k
in-domain examples ($0.837$), beats an SBERT prototype-cosine
baseline ($0.793$;~\citealp{reimers2019sbert}), and trails a $355$M
RoBERTa-large-MNLI zero-shot entailment classifier
($0.912$;~\citealp{liu2019roberta}) by $0.08$---all from nine
sentences of supervision and no labelled target-modality data.

\paragraph{The same direction across families.}
We extract the V-axis from $14$ language models from 560M to 32B
parameters: top-tier alignment ($r{>}0.85$ against a behavioural
valence reference) holds for all six Qwen3 sizes plus Mistral-7B;
Llama-4-Scout, Gemma-27B, and Gemma-4 sit in a middle tier; three
older small models (Pythia, TinyLlama, BLOOM) fall below detection
threshold (per-LLM table in Appendix~\ref{app:per-llm-deep-dive}).
Within the Qwen3 family, the V-axis is essentially scale-invariant
(off-diagonal $r{=}+0.995$ across $15$ pairs); across families,
correlations are looser ($r{=}+0.585$ over $48$ pairs). The
direction is a property of \emph{modern} LM training rather than of
parameter count.

\paragraph{Specificity.}
The recipe is concept-generic ($20$ further concepts produce a
working axis, $17$ exceed AUC $0.95$;
Appendix~\ref{app:concept-library}). Nonce-word ablation drops SST-2
to chance and random-Gaussian directions give chance lexicon recovery,
so the V-axis is a real semantic direction, not a noise artefact
(Appendix~\ref{app:specificity-controls}). An Arditi-style same-model
ablation~\citep{arditi2024refusal} on Qwen3-4B quantifies how much
sentiment lives along the V-axis: a logistic-regression probe on the
full features reaches SST-2 AUC $0.909$, drops to $0.907$ after
projecting the V-axis out and retraining ($z{=}{+}9.8$ above a
$20$-direction random-direction null). The V-axis carries detectable
sentiment signal but is not the sole sentiment-bearing direction.
